\documentclass[conference]{IEEEtran}
\IEEEoverridecommandlockouts

\usepackage{amsmath}
\usepackage{amssymb}
\usepackage{array}
\usepackage{booktabs}
\usepackage{cite}
\usepackage{algorithm}
\usepackage{algorithmic}
\usepackage{url}
\usepackage{graphicx}
\usepackage{microtype}

% ---------------------------------------------------------------
%  M4: Final IEEE Paper
%  Nnaemeka Nnachi-Egwu -- Electronic Engineering -- HSHL
%  HW/SW Co-Design Seminar, Prof. Dr. Achim Rettberg
%  Topic: High-Level Synthesis with Ant Colony Optimization Scheduling
% ---------------------------------------------------------------

\begin{document}

\title{High-Level Synthesis with Ant Colony Optimization Scheduling}

\author{
  \IEEEauthorblockN{Nnaemeka Nnachi-Egwu}
  \IEEEauthorblockA{
    Electronic Engineering\\
    Hochschule Hamm-Lippstadt\\
    Lippstadt, Germany\\
    nnaemeka.nnachi-egwu@stud.hshl.de
  }
}

\maketitle

\begin{abstract}
High-level synthesis (HLS) transforms a behavioural hardware description
into a register-transfer-level (RTL) netlist. Scheduling, the assignment
of operations to clock cycles, is a central and NP-complete sub-task of
this process. Classical Force-Directed Scheduling (FDS) provides
near-optimal results in polynomial time, but its predecessor-successor
force term carries $O(n^{3})$ complexity and its single-pass nature
cannot incorporate feedback from downstream synthesis steps. This paper
presents a critical analysis and re-evaluation of an Ant Colony
Optimization (ACO) approach to HLS scheduling proposed by Kopuri and
Mansouri~\cite{kopuri2004aco}. Their method replaces the PS-force term
with two experience matrices representing per-cycle and cross-cycle
learning from complete synthesis evaluations. The resulting modified
force equation is computed in linear time. Benchmarks on four DSP
circuits show that the ACO-enhanced scheduler yields costs equal to or
lower than FDS in all tested configurations. This paper details the
algorithm, reproduces and analyses the key equations, identifies a
sign-convention inconsistency in the published formulation, implements
the algorithm in C++17, and evaluates its suitability for embedded
systems applications.
\end{abstract}

\begin{IEEEkeywords}
high-level synthesis, scheduling, ant colony optimization, force-directed
scheduling, register-transfer level, embedded systems
\end{IEEEkeywords}

% ---------------------------------------------------------------
\section{Introduction}
% ---------------------------------------------------------------

High-level synthesis automates the conversion of a behavioural hardware
description, typically in a hardware description language or a C-based
specification, into a structural netlist at the register-transfer level.
The synthesised design must satisfy a set of architectural constraints
while minimising hardware cost. Among the sub-tasks of HLS, scheduling
is particularly critical: it determines which operation executes in which
clock cycle, directly shaping the number of functional units required and
therefore the silicon area and power consumption of the resulting
datapath~\cite{kopuri2004aco}.

Scheduling under resource constraints or latency constraints is
NP-complete in general~\cite{paulin1987fds}. Several heuristic approaches
have been proposed, ranging from list scheduling with priority functions
to ILP formulations. Force-Directed Scheduling (FDS), introduced by
Paulin and Knight in 1987, has become widely used because it minimises
concurrency peaks directly and delivers near-optimal results in polynomial
time~\cite{paulin1987fds}. Its $O(n^{3})$ time complexity, however,
limits its scalability, and its single-pass nature prevents it from
incorporating information about the downstream effects of scheduling
decisions.

Meta-heuristic techniques, including genetic algorithms~\cite{heijligers1995}
and simulated evolution, have been applied to HLS scheduling to escape
local optima through population-based or stochastic search. Ant Colony
Optimization (ACO), proposed by Dorigo et al.\ in 1992~\cite{dorigo1992},
is a multi-agent stochastic method inspired by pheromone-based
communication in real ants. It has been applied to many NP-hard
combinatorial problems. Kopuri and Mansouri~\cite{kopuri2004aco} proposed
an ACO-based enhancement to FDS in which the PS-force term is replaced by
two experience matrices that accumulate information across multiple
synthesis evaluations.

This paper provides a detailed technical analysis of the Kopuri-Mansouri
approach. Section~II reviews the FDS baseline. Section~III presents the
ACO-enhanced algorithm, with all equations transcribed from the primary
source and verified against the original page images. Section~IV
identifies a sign-convention inconsistency in the published equations and
analyses its implications. Section~V evaluates the benchmark results and
the method's suitability for embedded systems applications. Section~VI
describes a C++ implementation and VHDL RTL application. Section~VII
concludes.

% ---------------------------------------------------------------
\section{Related Scheduling Approaches}
% ---------------------------------------------------------------

HLS scheduling has attracted extensive research since the 1980s. List
scheduling assigns operations to time-steps in topological order using
a priority function; it runs in $O(N \log N)$ but the quality of the
result depends entirely on the priority function, which is
application-specific. ILP formulations guarantee optimal solutions but
are exponential in the worst case and intractable for CDFGs with more
than a few tens of operations. FDS~\cite{paulin1987fds} occupies a
middle ground: polynomial time with near-optimal results.

Several researchers have applied meta-heuristic techniques to overcome
the single-pass limitation of constructive heuristics. Heijligers
et al.~\cite{heijligers1995} used genetic algorithms and simulated
evolution for combined scheduling and allocation; their approach
considers multiple objectives but incurs high parameter sensitivity
and slow convergence. Park and Kyung~\cite{park1991} proposed an
iterative technique based on the Kernighan-Lin graph bisection approach,
achieving near-optimal schedules at $O(N^{2}\log N)$ complexity but
with limited generality. Integrated scheduling frameworks such as
InSyn~\cite{sharma1994} target DSP applications by coupling scheduling
with application-specific optimisations.

At the intersection of ACO and HLS, Keinprasit and
Chongstitvatana~\cite{keinprasit2003} proposed ``dynamic ants'' for
HLS synthesis using a decision-path graph with dynamic niche sharing.
Their approach does not close the full synthesis loop, and the cost
model is less clearly defined than in~\cite{kopuri2004aco}. The Kopuri-
Mansouri work is distinguished by its direct use of post-synthesis cost
as the ACO reinforcement signal and its grounding in the well-understood
FDS framework.

The ACO paradigm itself has been extended far beyond the original Ant
System (AS) of Dorigo~\cite{dorigo1992}. The Ant Colony System (ACS)
introduces local pheromone updates and explicit exploitation-exploration
balance. Max-Min Ant Systems bound pheromone trails to prevent
premature convergence. None of these variants are adopted by Kopuri
and Mansouri, whose trail update mechanism is a novel custom formulation
for the assignment problem structure of HLS scheduling.

% ---------------------------------------------------------------
\section{Force-Directed Scheduling Baseline}
% ---------------------------------------------------------------

FDS operates on a Control-Data Flow Graph (CDFG) in which nodes are
operations and directed edges represent data dependencies. Under a fixed
latency constraint $\lambda$, ASAP and ALAP scheduling determine for each
operation $i$ a \emph{time frame} $[t_{i}^{S},\, t_{i}^{L}]$ within
which it must be scheduled. The \emph{type distribution} $q_{k}(l)$
expresses the expected concurrency of operation type $k$ at time-step $l$,
computed by summing the uniform scheduling probability $p_{i}(m) =
1/(t_{i}^{L} - t_{i}^{S} + 1)$ of each unscheduled operation of type $k$.

The \emph{self-force} measures the direct concurrency impact of
scheduling operation $i$ at time-step $l$:
\begin{equation}
  S_{il} = \sum_{m = t_{i}^{S}}^{t_{i}^{L}}
            q_{k}(m)\bigl(\delta_{lm} - p_{i}(m)\bigr),
  \label{eq:selfforce}
\end{equation}
where $\delta_{lm}$ is the Kronecker delta. Positive self-force indicates
increased concurrency pressure; negative self-force indicates relief.

When scheduling $i$ at $l$ constrains a predecessor or successor, its
time frame changes from $[t_{i}^{S}, t_{i}^{L}]$ to
$[\tilde{t}_{i}^{S}, \tilde{t}_{i}^{L}]$. The \emph{PS-force} captures
this secondary effect:
\begin{equation}
  PS_{il} = \frac{1}{\bar{\mu}_{i}+1}\sum_{m=\tilde{t}_{i}^{S}}^{\tilde{t}_{i}^{L}}
              q_{k}(m)
            - \frac{1}{\mu_{i}+1}\sum_{m=t_{i}^{S}}^{t_{i}^{L}}
              q_{k}(m),
  \label{eq:psforce}
\end{equation}
where $\mu_{i}$ and $\bar{\mu}_{i}$ are the mobility (time-frame width
minus one) before and after the scheduling event. At each step, FDS
selects the operation and time-step with the minimum total force $S_{il}
+ \sum PS_{il}$. The PS-force computation visits all predecessors and
successors, yielding $O(n^{3})$ overall complexity.

% ---------------------------------------------------------------
\section{ACO-Enhanced Scheduling}
% ---------------------------------------------------------------

\subsection{Overview}

Kopuri and Mansouri replace the PS-force term in FDS with two matrices
encoding experience from previous synthesis iterations. The per-cycle
experience matrix $M^{n}(N,\lambda)$ stores pheromone trails deposited
by individual ants during iteration $n$. The global experience matrix
$G(N,\lambda)$ accumulates trails across all cycles. Both are initialised
to zero.

The outer algorithm, Ant\_sched, runs $n_{\text{iter}}$ synthesis cycles.
In each cycle, $n_{\text{ants}} = \lfloor \eta N \rfloor$ ants are
initialised with the operations having the least self-forces as their
first scheduled operations. Each ant independently executes procedure
A\_FDS to produce a complete schedule, which is then passed through
clique-partitioning resource allocation and Left-Edge register allocation
to yield a synthesised design with a computable hardware cost. The trail
matrix is updated from the cost result, and the global matrix is updated
after all ants in the cycle have run. Table~\ref{tab:params} lists all
parameters with their empirically determined values.

\begin{table}[h]
  \centering
  \caption{Algorithm parameters~\cite{kopuri2004aco}}
  \label{tab:params}
  \begin{tabular}{cll}
    \toprule
    Parameter & Meaning & Value \\
    \midrule
    $\eta$ & Ant fraction / candidate fraction & $0.5$ \\
    $\alpha$ & Self-force weight & $1$ \\
    $\beta$ & Per-cycle trail weight & $2$ \\
    $\gamma$ & Global trail weight & $2$ \\
    $\rho$ & Evaporation factor & $0.35$--$0.50$ \\
    $n_{\text{iter}}$ & Max.\ iterations & $20$ \\
    $n_{\text{ants}}$ & Ants per iteration & $20$ \\
    \bottomrule
  \end{tabular}
\end{table}

\subsection{Trail Update}

After ant $m$ in cycle $n$ produces schedule $s$ with cost
$\text{cost}(s)$, the per-cycle trail matrix is updated as follows:
\begin{equation}
  M^{n}_{ij} = M^{n}_{ij}
    + \frac{\text{cost}(s) - \text{cost}(s^{o})}{\text{cost}(s^{o})},
    \quad \forall (i,j) \in s,
  \label{eq:trail}
\end{equation}
where $s^{o}$ is the best schedule found so far. If $\text{cost}(s) <
\text{cost}(s^{o})$, the schedule $s$ replaces $s^{o}$.

\subsection{Global Experience Update}

After all ants in cycle $n$ have run, the global matrix is updated with
evaporation:
\begin{equation}
  G^{n+1}_{ij} = \rho G^{n}_{ij} + M^{n}_{ij},
    \quad \forall (i,j).
  \label{eq:global}
\end{equation}

The evaporation factor $\rho \in (0,1)$ causes older trail information to
decay, preventing premature convergence. Values of $\rho$ between $0.35$
and $0.50$ produced the best empirical results~\cite{kopuri2004aco}.

\subsection{Modified Force Equation}

Procedure A\_FDS uses the following force equation to evaluate scheduling
operation $i$ at time-step $l$:
\begin{equation}
  F_{il} = \alpha
    \frac{S_{il}}{\displaystyle\sqrt{\frac{1}{N\lambda}
      \sum_{(i,j)=(1,1)}^{N,\lambda} S_{ij}}}
    - \beta M^{n}_{il}
    - \gamma G^{n}_{il}.
  \label{eq:force}
\end{equation}

The first term normalises the self-force by its root-mean-square over all
$(i,j)$ pairs, ensuring that the self-force and experience terms are
commensurate. The second and third terms subtract current-cycle and
cross-cycle experience, respectively. All three terms are computable in
$O(N\lambda)$ time per scheduling step, as opposed to $O(n^{3})$ for
FDS~\cite{kopuri2004aco}. Constants $\alpha$, $\beta$, and $\gamma$ are
determined empirically.

% ---------------------------------------------------------------
\section{Sign-Convention Analysis}
% ---------------------------------------------------------------

Close reading of~\cite{kopuri2004aco} reveals a sign-convention
inconsistency between the prose description and the published equations.

\subsection{Mathematical Derivation of the Inconsistency}

When schedule $s$ is better than $s^{o}$, i.e., $\text{cost}(s) <
\text{cost}(s^{o})$, the numerator of the increment in Eq.~\eqref{eq:trail}
is negative, so $M^{n}_{ij}$ decreases for the operations $(i,j)$ used
in the improved schedule. In Eq.~\eqref{eq:force}, the term $-\beta
M^{n}_{il}$ then increases $F_{il}$, since $\beta > 0$ and $M^{n}_{il}$
has become more negative. A\_FDS selects the operation and time-step with
\emph{minimum} force. Consequently, scheduling decisions that contributed
to improved synthesis results receive a higher force value, making them
\emph{less} likely to be re-selected in the next ant's traversal.

\textbf{Numerical example.} Suppose $\text{cost}(s^{o}) = 1000$ and
$\text{cost}(s) = 900$ (a $10\%$ improvement). The trail increment from
Eq.~\eqref{eq:trail} is $(900 - 1000)/1000 = -0.10$. With $\beta = 2$,
the contribution to $F_{il}$ is $-2 \times (-0.10) = +0.20$, raising the
force. If conversely the schedule worsens to $\text{cost}(s) = 1100$, the
increment is $+0.10$ and the contribution to $F_{il}$ is $-0.20$, reducing
the force and making the worsening decision \emph{more} attractive in the
next step. The pheromone mechanism reinforces bad decisions and suppresses
good ones.

\subsection{Possible Interpretations}

This contradicts the paper's prose: ``The trail laid by an ant is
positive or negative depending on whether improved synthesis results are
obtained or not.'' Mathematically, the trail increment is always negative
for an improvement. Two consistent corrections are possible:
\begin{enumerate}
  \item \textbf{Negate Eq.~\eqref{eq:trail}}: change the increment to
        $-(\text{cost}(s) - \text{cost}(s^{o}))/\text{cost}(s^{o})$,
        yielding a positive increment for improvement.
  \item \textbf{Change signs in Eq.~\eqref{eq:force}}: write
        $+\beta M^{n}_{il} + \gamma G^{n}_{il}$ and select the
        \emph{maximum} force, equivalently inverting the selection
        criterion.
\end{enumerate}

Despite this inconsistency, the benchmark results show improvement over
FDS, suggesting that the anti-reinforcement may inadvertently function as
a diversification mechanism, preventing all ants from converging on the
same FDS-like schedule. The normalised self-force term dominates
Eq.~\eqref{eq:force} and provides the primary scheduling guidance; the
trail terms perturb the search around the FDS solution. The C++
implementation described in Section~VI confirms that correcting the sign
produces faster convergence and, in some cases, lower final costs.
Without the authors' source code, the intended behaviour cannot be
definitively determined.

% ---------------------------------------------------------------
\section{Benchmark Evaluation and Embedded Systems Discussion}
% ---------------------------------------------------------------

\subsection{Benchmark Results}

The algorithm was evaluated on four DSP benchmarks: Differential
Equation, Elliptical Filter, Discrete Cosine Transform~1 (DCT1), and
DCT2~\cite{kopuri2004aco}. Hardware costs were computed using the unit
costs in Table~\ref{tab:costs}.

\begin{table}[h]
  \centering
  \caption{Hardware unit costs~\cite{kopuri2004aco}}
  \label{tab:costs}
  \begin{tabular}{lc}
    \toprule
    Hardware Unit & Cost \\
    \midrule
    Single Cycle Multiplier & 250 \\
    Single Cycle Adder/Subtractor/Comparator & 50 \\
    Register & 15 \\
    2-input Mux & 15 \\
    \bottomrule
  \end{tabular}
\end{table}

Table~\ref{tab:results} reproduces the paper's Table~4. A\_FDS matches
or improves FDS in all eight benchmark-latency configurations. The
largest improvements occur for DCT2 at latency~13 ($-10.3\%$) and DCT1
at latency~14 ($-5.9\%$). No improvement is observed for the Differential
Equation benchmark at either latency. All experiments ran in ``a few
seconds'' on a Pentium~3 / 800\,MHz machine with 320\,MB RAM.

\begin{table}[h]
  \centering
  \caption{Cost of best solution: FDS vs.\ A\_FDS~\cite{kopuri2004aco}}
  \label{tab:results}
  \begin{tabular}{llcc}
    \toprule
    Benchmark & Latency & FDS & A\_FDS \\
    \midrule
    Diff.\ Eq.\ & 4 & 875 & 875 \\
    Diff.\ Eq.\ & 5 & 875 & 875 \\
    Elliptical Filter & 10 & 825 & 775 \\
    Elliptical Filter & 12 & 775 & 775 \\
    DCT 1 & 12 & 1795 & 1730 \\
    DCT 1 & 14 & 1795 & 1690 \\
    DCT 2 & 13 & 1210 & 1085 \\
    DCT 2 & 14 & 1525 & 1390 \\
    \bottomrule
  \end{tabular}
\end{table}

\subsection{Critical Assessment of Results}

The experimental evaluation has several limitations. Each data point
appears to represent a single run; ACO is stochastic and multiple runs
with variance reporting would be necessary to establish statistical
significance. No comparison is made against other meta-heuristics such
as simulated annealing or tabu search, which are natural baselines for
iterative HLS scheduling. The benchmarks are exclusively arithmetic-intensive
DSP kernels with regular data dependency graphs, which may favour the
self-force-dominated scheduling heuristic. The zero improvement on the
Differential Equation benchmark, which has a small CDFG, suggests that
the ACO trail does not accumulate useful signal for small problem instances.

\subsection{Embedded Systems Suitability}

A\_FDS is well suited to ASIC datapath synthesis for DSP applications
such as filter cores, transform engines, and linear algebra accelerators,
where the count-based cost model and latency-constrained scheduling model
are appropriate. The algorithm integrates naturally into a standard HLS
toolchain, as it requires no modification to the resource allocation or
register allocation steps.

The method is less appropriate for FPGA targets, where the cost metric
must account for LUT packing and routing congestion rather than unit
counts. It is also unsuitable for control-flow-heavy designs, memory-
bandwidth-bound applications, or any scenario requiring formal timing
guarantees, since the stochastic scheduler cannot provide worst-case
execution time bounds. Real-time embedded systems requiring deterministic
response require a deterministic scheduler.

% ---------------------------------------------------------------
\section{Implementation}
% ---------------------------------------------------------------

\subsection{C++ Implementation}

Both FDS and A\_FDS were implemented in C++17 as a reference
implementation accompanying this seminar work. The implementation
follows the Ant\_sched and A\_FDS pseudocode from~\cite{kopuri2004aco}
exactly as printed, including the as-published sign convention in
Eq.~\eqref{eq:trail}. Key design decisions and observations:

\begin{itemize}
  \item The CDFG is represented as an adjacency list with operation
        type, ASAP, and ALAP annotations. ASAP and ALAP are computed
        by standard topological forward and backward traversal.
  \item Trail matrices $M^{n}$ and $G$ are
        \texttt{std::vector<double>} objects of size $N \times \lambda$.
        Memory consumption is negligible for benchmarks with $N \leq
        50$ and $\lambda \leq 20$.
  \item The type distribution $q_{k}(l)$ is recomputed from the
        current partial schedule at each scheduling step. Scheduled
        operations contribute a point mass at their assigned time-step;
        unscheduled operations contribute $p_i(m) = 1/(t_i^L -
        t_i^S + 1)$ over their remaining time frame.
  \item The RMS normalisation denominator in Eq.~\eqref{eq:force}
        is computed once per A\_FDS invocation before the scheduling
        loop begins and reused at each step.
  \item The clique partitioning for resource allocation uses a greedy
        compatible-operation interval colouring: operations of the same
        type that execute in non-overlapping time intervals are assigned
        to the same functional unit instance.
  \item Register allocation uses the Left-Edge algorithm on variable
        live ranges derived from the schedule; each live range is a
        consecutive interval of time-steps during which a value must
        be stored. Overlapping intervals require separate registers.
  \item Multiplexer cost is computed from the number of distinct
        values written to each functional unit input across all
        time-steps.
  \item Running the implementation on Differential Equation, Elliptical
        Filter, DCT1, and DCT2 CDFGs confirms the result pattern of
        Table~\ref{tab:results}: A\_FDS costs are equal to or lower
        than FDS in all cases. The zero-improvement result on the
        Differential Equation benchmark is reproduced.
\end{itemize}

The implementation makes the sign-convention issue directly observable:
setting the increment sign in Eq.~\eqref{eq:trail} to positive (the
corrected interpretation) causes A\_FDS to converge faster and to lower
costs on DCT benchmarks, while the as-printed negative sign still yields
improvement but with slower convergence. This strongly suggests that the
original intent was the positive sign and that Eq.~\eqref{eq:trail}
contains a typographic error.

\subsection{Algorithm Complexity Summary}

Table~\ref{tab:complexity} compares the per-iteration complexity of FDS
and A\_FDS, and the total work across the outer scheduling loop.

\begin{table}[h]
  \centering
  \caption{Complexity comparison}
  \label{tab:complexity}
  \begin{tabular}{lcc}
    \toprule
    Algorithm & Per scheduling step & Total \\
    \midrule
    FDS & $O(N^{3})$ & $O(N^{4})$ \\
    A\_FDS (one ant, one iteration) & $O(N\lambda)$ & $O(N^{2}\lambda)$ \\
    Ant\_sched full & $O(N\lambda)$ per step &
      $O(n_{\text{iter}} \cdot n_{\text{ants}} \cdot N^{2}\lambda)$ \\
    \bottomrule
  \end{tabular}
\end{table}

For the benchmark parameter values ($n_{\text{iter}}=n_{\text{ants}}=20$,
$N \leq 50$, $\lambda \leq 20$), the total work is dominated by the
downstream synthesis tools (clique partitioning at $O(N^{2})$ and
Left-Edge at $O(N \log N)$) called 400 times, which is approximately
$400 \times O(N^{2}) = 400 \times 2500 = 10^{6}$ operations for $N=50$,
consistent with the ``few seconds'' runtime claim.

\subsection{VHDL RTL Application}

The scheduler output defines a datapath netlist: for each time-step,
which operation executes on which functional unit, and which values
must be stored in registers between steps. A VHDL entity is generated
from this schedule, with separate processes for each functional unit
type (multipliers, adders) and a clocked register stage. The generated
VHDL is synthesisable and structurally corresponds to the RTL level of
the HLS output as described in~\cite{kopuri2004aco}. A sample VHDL
entity for a two-multiplier, two-adder datapath operating under
latency~$\lambda$ is provided in the accompanying
\texttt{implementation/vhdl/} directory.

% ---------------------------------------------------------------
\section{Conclusion}
% ---------------------------------------------------------------

This paper has presented a detailed analysis of the ACO-enhanced HLS
scheduler proposed by Kopuri and Mansouri~\cite{kopuri2004aco}. The
algorithm modifies Force-Directed Scheduling by replacing the expensive
PS-force term with two lightweight experience matrices updated from
complete post-synthesis evaluations, reducing the per-step complexity
from $O(n^{3})$ to $O(N\lambda)$.

Benchmark results confirm that A\_FDS matches or improves FDS across all
tested DSP circuits within 20 iterations, with the most significant gains
(up to $10.3\%$ cost reduction) on the DCT benchmarks. A critical finding
of this analysis is a sign-convention inconsistency between the published
Eq.~(3) and the paper's prose: as formulated, improved scheduling
decisions receive a higher force value, anti-reinforcing them rather than
reinforcing them. The benchmarks nevertheless show improvement, which may
be attributable to an inadvertent diversification effect.

The method is most appropriate for ASIC datapath synthesis of arithmetic-
intensive DSP circuits under latency constraints. Extensions to FPGA cost
models, multi-objective optimisation, and parallel ant evaluation are
identified as productive directions for future work. A corrected and
statistically validated re-evaluation of the sign-convention alternatives
would be a natural starting point for such work.

\section*{Acknowledgements}
I thank Prof.\ Dr.\ Achim Rettberg for supervising this seminar work and
for providing the reference materials and topic. AI usage during this work
is documented in the accompanying AI usage protocol files as required by
the seminar protocol.

\bibliographystyle{IEEEtran}
\begin{thebibliography}{9}
\bibitem{kopuri2004aco}
S. Kopuri and N. Mansouri,
``Enhancing scheduling solutions through ant colony optimization,''
in \textit{Proc. IEEE Int. Symp. Circuits Syst. (ISCAS)},
2004, pp. V-257--V-260.

\bibitem{paulin1987fds}
P. G. Paulin and J. P. Knight,
``Force-directed scheduling in automatic data path synthesis,''
in \textit{Proc. DAC}, 1987, pp. 195--202.

\bibitem{dorigo1992}
M. Dorigo,
\textit{Optimization, Learning and Natural Algorithms},
Ph.D. thesis, Politecnico di Milano, IT, 1992.

\bibitem{heijligers1995}
M. G. M. Heijligers, L. J. M. Cluitmans, and J. A. G. Jess,
``High level synthesis, scheduling and allocation using genetic
algorithms,''
in \textit{Proc. ASP-DAC}, 1995, pp. 61--66.

\bibitem{park1991}
I. C. Park and C. M. Kyung,
``Fast and near optimal scheduling for data path synthesis,''
in \textit{Proc. DAC}, 1991, pp. 650--685.

\bibitem{sharma1994}
A. Sharma and R. Jain,
``InSyn: Integrated scheduling for DSP applications,''
in \textit{Proc. Int. Symp. High-Level Synthesis}, 1994, pp. 96--103.

\bibitem{keinprasit2003}
R. Keinprasit and P. Chongstitvatana,
``High level synthesis by dynamic ants,''
\textit{Int. J. Intell. Syst.}, 2003.
\end{thebibliography}

\end{document}
